Consider a spin Young diagram with two rows containing $N_1$ and $N_2$ boxes, respectively ($N_1+N_2=N$, $N_1\geq N_2$). Each of the first $N_2$ columns has two boxes, so the spinor must be antisymmetric in each corresponding pair of indices. By contrast, it must be symmetric in the indices occupying the final $n=N_1-N_2$ boxes of the first row. As we know, a rank-$N$ spinor of this kind reduces to a symmetric spinor of rank $n$, and the latter has total spin $S=n/2$. In terms of the Young diagrams for the coordinate functions, therefore, a diagram having $n$ one-box rows corresponds to a state whose total spin is $S=n/2$. For even $N$, the total spin may take the integer values from 0 through $N/2$; for odd $N$, it may take the half-integer values from $1/2$ through $N/2$, as required.
